\subsection*{Propositions}

\begin{remark*}[Source note]
The following construction propositions are Euclid's Propositions I.1--I.3 in
modern theorem form, following the Heath translation of the
\textit{Elements} \citep[Book I, Propositions 1--3]{euclid-elements-heath}.
\end{remark*}


\begin{theorem}[Euclid I.1: Constructing an equilateral triangle]
\label{thm:euclid-i-1}
Given a finite straight line segment $AB$, there exists an equilateral
triangle $ABC$ constructed on $AB$.
\end{theorem}
\begin{remark*}[Proof]
\hyperref[prf:euclid-i-1]{Go to proof.}
\end{remark*}
\begin{remark*}[Standard quantified statement]
This is a classical Euclidean source-text clause. In this chapter it is treated
as a named primitive statement rather than expanded into a modern first-order
formalization.
\end{remark*}
\begin{lra-not-visible}
\begin{dependencies}
\begin{itemize}
  \item \hyperref[def:euclid-point]{Point}
  \item \hyperref[def:euclid-straight-line]{Straight line}
  \item \hyperref[def:euclid-circle]{Circle}
  \item \hyperref[def:euclid-triangle-side-classes]{Equilateral triangle}
  \item \hyperref[ax:euclid-postulate-draw-straight-line]{Euclid Postulate I}
  \item \hyperref[ax:euclid-postulate-describe-circle]{Euclid Postulate III}
  \item \hyperref[ax:euclid-common-notion-things-equal-to-same]{Euclid Common Notion I}
\end{itemize}
\end{dependencies}
\end{lra-not-visible}

\begin{theorem}[Euclid I.2: Copying a segment from a given point]
\label{thm:euclid-i-2}
Given a point $A$ and a finite straight line segment $BC$, there exists a
straight line segment beginning at $A$ and equal to $BC$.
\end{theorem}
\begin{remark*}[Proof]
\hyperref[prf:euclid-i-2]{Go to proof.}
\end{remark*}
\begin{remark*}[Standard quantified statement]
This is a classical Euclidean source-text clause. In this chapter it is treated
as a named primitive statement rather than expanded into a modern first-order
formalization.
\end{remark*}
\begin{lra-not-visible}
\begin{dependencies}
\begin{itemize}
  \item \hyperref[def:euclid-point]{Point}
  \item \hyperref[def:euclid-straight-line]{Straight line}
  \item \hyperref[def:euclid-circle]{Circle}
  \item \hyperref[thm:euclid-i-1]{Euclid I.1}
  \item \hyperref[ax:euclid-postulate-draw-straight-line]{Euclid Postulate I}
  \item \hyperref[ax:euclid-postulate-produce-straight-line]{Euclid Postulate II}
  \item \hyperref[ax:euclid-postulate-describe-circle]{Euclid Postulate III}
  \item \hyperref[ax:euclid-common-notion-things-equal-to-same]{Euclid Common Notion I}
\end{itemize}
\end{dependencies}
\end{lra-not-visible}

\begin{theorem}[Euclid I.3: Cutting off an equal segment]
\label{thm:euclid-i-3}
Given two unequal finite straight line segments, there exists a construction
which cuts off from the greater a straight line segment equal to the less.
\end{theorem}
\begin{remark*}[Proof]
\hyperref[prf:euclid-i-3]{Go to proof.}
\end{remark*}
\begin{remark*}[Standard quantified statement]
This is a classical Euclidean source-text clause. In this chapter it is treated
as a named primitive statement rather than expanded into a modern first-order
formalization.
\end{remark*}
\begin{lra-not-visible}
\begin{dependencies}
\begin{itemize}
  \item \hyperref[def:euclid-straight-line]{Straight line}
  \item \hyperref[def:euclid-circle]{Circle}
  \item \hyperref[thm:euclid-i-2]{Euclid I.2}
  \item \hyperref[ax:euclid-postulate-describe-circle]{Euclid Postulate III}
  \item \hyperref[ax:euclid-common-notion-things-equal-to-same]{Euclid Common Notion I}
\end{itemize}
\end{dependencies}
\end{lra-not-visible}
